\documentclass[12pt]{article}
\usepackage{fullpage}
\usepackage{amsmath,amssymb,amsthm,mathtools,array}
\usepackage[hidelinks]{hyperref}

\newcommand{\Z}{\mathbb Z}
\newcommand{\Q}{\mathbb Q}
\newcommand{\C}{\mathbb C}
\newcommand{\PP}{\mathbb P}
\newcommand{\F}{\mathbb F}
\newcommand{\NS}{\operatorname{NS}}
\newcommand{\Pic}{\operatorname{Pic}}
\newcommand{\rk}{\operatorname{rk}}
\newcommand{\im}{\operatorname{im}}
\newcommand{\Bl}{\operatorname{Bl}}
\newcommand{\calX}{\mathcal X}
\newcommand{\calW}{\mathcal W}
\newcommand{\calZ}{\mathcal Z}
\newcommand{\calL}{\mathcal L}
\newcommand{\II}{\mathrm{II}}
\newcommand{\I}{\mathrm{I}}
\newcommand{\LK}{\Lambda_{\mathrm{K3}}}
\newtheorem{theorem}{Theorem}[section]
\newtheorem{proposition}[theorem]{Proposition}
\newtheorem{lemma}[theorem]{Lemma}
\newtheorem{corollary}[theorem]{Corollary}
\theoremstyle{definition}
\newtheorem{definition}[theorem]{Definition}
\theoremstyle{remark}
\newtheorem{remark}[theorem]{Remark}

\begin{document}

\begin{center}
\textbf{Mirror fibrations for very general \(H\)-polarised Tyurin
degenerations of K3 surfaces}
\end{center}

\section{Introduction}

\subsection{The statement}

Throughout, \(H\) is the hyperbolic plane lattice (even, unimodular,
of signature \((1,1)\)), \(E_8\) the negative definite \(E_8\) root
lattice, \(\LK=H^{\oplus3}\oplus E_8^{\oplus2}\) the K3 lattice, and
\(M=H\oplus E_8\oplus E_8\). For an even lattice \(N\) of signature
\((1,\rk N-1)\), an \emph{\(N\)-quasipolarised} K3 surface is a K3
surface \(Y\) with a primitive embedding \(N\hookrightarrow\NS(Y)\)
satisfying the chamber convention of \cite[Definition~3.19]{GT} (the
embedding sends a fixed very irrational positive vector of
\(N\otimes\mathbb R\) to a nef and big class; the convention goes back
to \cite{Do}); in particular the image of \(N\) contains a nef and big
class. For very general such \(Y\) (outside countably
many Noether--Lefschetz divisors of the moduli space)
\(N\simeq\NS(Y)\). In Dolgachev's mirror symmetry for
lattice-polarised K3 surfaces \cite{Do} the lattices \(H\) and \(M\)
are mirror to each other, as \(H^\perp_{\LK}\simeq H\oplus M\).

A \emph{Kulikov model} is a proper flat family
\(\pi\colon\calX\to\Delta\) over a disc with smooth total space,
trivial canonical bundle, reduced normal crossings central fibre and
smooth K3 general fibre \cite{Ku}. A \emph{Tyurin degeneration} is a
Kulikov model whose central fibre \(X_0=V_1\cup V_2\) is the union of
two smooth rational surfaces meeting transversally along a smooth
elliptic curve \(C\) (a Kulikov model of \emph{Type~II}: the double
curves are elliptic). By adjunction \(C\) is anticanonical on each
\(V_i\), and smoothness of \(\calX\) makes \(X_0\) \(d\)-semistable
\cite{Fr83}: \(N_{C/V_1}\otimes N_{C/V_2}\simeq\mathcal O_C\), so
\((C^2)_{V_1}+(C^2)_{V_2}=0\).

Following the proposal of Doran, Harder and Thompson \cite{DHT} that
Tyurin degenerations should be mirror to fibrations over \(\PP^1\),
Giovenzana and Thompson \cite{GT} define when an \(N\)-polarised
Tyurin degeneration (Definition~\ref{def:Hpol} below, for \(N=H\)) and
an elliptic fibration on the mirror K3 surface, split into two pieces
by a loop in \(\PP^1\), form a \emph{mirror pair}
\cite[Definition~5.1]{GT}; the relevant notions are recalled in
\S\ref{sec:recall}.

\begin{definition}\label{def:Hpol}
A Tyurin degeneration \(\pi\colon\calX\to\Delta\) is
\emph{\(H\)-polarised} if
(1) (monodromy invariance) over the punctured disc \(\Delta^*\)
there is a trivial local subsystem \(\calL\subset R^2\pi_*\Z\) with
fibre \(H\), such that for each \(t\in\Delta^*\) the fibre
\(\calL_t\subset H^2(X_t;\Z)\) defines an \(H\)-quasipolarisation of
\(X_t\); and
(2) (positivity) there is an effective, flat (over \(\Delta\), i.e.\
with no component inside \(X_0\)) and nef Cartier divisor
\(A\) on \(\calX\) whose restriction \(A_t\) to every fibre \(X_t\),
\(t\in\Delta^*\), is nef and big with \([A_t]\in\calL_t\), and whose
restriction to each component satisfies \((A|_{V_i})^2>0\),
\(i=1,2\).
\end{definition}

This is \cite[Definition~3.20]{GT} with two strengthenings of (2)
that are part of our hypotheses: the membership \([A_t]\in\calL_t\),
which corrects the arXiv version of that definition, and the
componentwise positivity \((A|_{V_i})^2>0\).

\begin{theorem}\label{thm:main}
Let \(\pi\colon\calX\to\Delta\) be an \(H\)-polarised Tyurin
degeneration of K3 surfaces with central fibre \(X_0=V_1\cup V_2\).
Assume it is very general, in the sense that the image of
\(0\in\Delta\) under the map from \(\Delta\) to the
Alexeev--Brunyate--Engel compactification \cite{ABE} of the moduli
space of \(H\)-polarised K3 surfaces lies in no Noether--Lefschetz
locus. Let \(Y\) be a very general \(M\)-quasipolarised K3 surface.
Then there exist an elliptic fibration \(Y\to\PP^1\) and an allowable
loop \(\gamma\subset\PP^1\) \cite[Definition~4.11]{GT} such that the
\(M\)-quasipolarisation is compatible with the fibration and
\(\gamma\) \cite[Definition~4.21]{GT}, and such that \(\calX\to\Delta\)
and \((Y\to\PP^1,\gamma)\) form a mirror pair in the sense of
\cite[Definition~5.1]{GT}.
\end{theorem}

\subsection{Conventions and interpretation}\label{sec:conventions}

We use the definitions of \cite{GT} in arXiv version~3.
Noether--Lefschetz loci on the compactification are the closures of
the Noether--Lefschetz loci of its interior. All lattices and homology
groups are integral; \(x\cdot y\) or \((x,y)\) denotes the bilinear
form, \(x^2=(x,x)\), and a \emph{root} is a vector of square \(-2\).

The degeneration determines more than the image point of \(0\) in the
ABE compactification: it determines the \emph{marked} (``mixed-marked''
in \cite{AE}) Type~II period of its central fibre, namely the isotropic
plane \(I\) of the monodromy (\S\ref{sec:degside}) and the restriction
homomorphism \(\varphi\) of \S\ref{sec:period}. We always work with
this marked period. The distinction matters: over one of the two
Type~II boundary strata the boundary normalisation map of
\cite[Theorem~7.58]{ABE} has degree eight, so a coarse boundary point
does not determine a marked period; we never pass from one to the
other, and no base change or level cover is introduced.

\subsection{Strategy}\label{sec:strategy}

Both sides produce an even unimodular lattice of signature \((1,17)\)
with a distinguished sublattice and two mutually orthogonal sublattices
(one for each component of \(X_0\), respectively each piece of the
split fibration); a mirror pair is essentially an isometry of these
rank-\(18\) lattices matching the embedded sublattices and lifting to
the K3 lattices (compare \cite[Proposition~5.3]{GT}). On the
degeneration side the hypotheses do not describe \(X_0\) explicitly.
Section~\ref{sec:period} shows that very generality forces
the kernel of the restriction homomorphism to be the polarising
lattice \(L\simeq H\), so the rank-\(18\) lattice is \(L\oplus R_X\)
with \(R_X\simeq E_8^2\) or \(D_{16}^+\) (notation of
\S\ref{sec:recall}). Section~\ref{sec:reference} constructs, for the
given marked period, an explicit Kulikov model \(\calW\) isomorphic to
\(\calX\) over \(\Delta^*\) and in codimension one. Since Kulikov
models are unique only up to elementary modifications (flops along
curves in the central fibre), \(\calW\) need not be \(\calX\);
Section~\ref{sec:models}
transports the divisor \(A\) to \(\calW\) and compares the contractions
of \(\calW\) and \(\calX\) defined by \(A\). Componentwise positivity
of \(A\) leaves exactly three central fibres, \(\mathrm E_0,\mathrm E_1\)
(for \(R_X=E_8^2\)) and \(\mathrm D\) (for \(R_X=D_{16}^+\)), whose
embedded lattices are computed in \S\ref{sec:latticesX}. On the
fibration side, \S\ref{sec:fibrations} exhibits two elliptic
fibrations on \(Y\) (the \emph{standard} one, with singular fibres
\(2\,\II^*+4\,\I_1\), and the \emph{alternate} one, with
\(\I_{12}^*+6\,\I_1\)) and three allowable splitting loops
\(\gamma_0,\gamma_1,\gamma_D\); \S\ref{sec:latticesY} proves
compatibility and computes the embedded lattices, which match the
degeneration cases according to
\(\mathrm E_0\leftrightarrow(\text{standard},\gamma_0)\),
\(\mathrm E_1\leftrightarrow(\text{standard},\gamma_1)\),
\(\mathrm D\leftrightarrow(\text{alternate},\gamma_D)\).
The delicate points are integral: the small piece of \(\gamma_D\) has
``degree \(8'\)'' rather than \(8\), and the large piece carries an
index-two parity condition (\S\ref{sec:latticesY}).
Section~\ref{sec:final} lifts the isometry to the K3 lattices and
checks every clause of \cite[Definition~5.1]{GT}.

\section{Notation and recollections}\label{sec:recall}

\subsection{Lattices}

\(D_{16}\) is the negative definite \(D_{16}\) root lattice and
\(D_{16}^+\) its even unimodular index-two overlattice; \(E_8^2\) and
\(D_{16}^+\) are the only even unimodular negative definite lattices of
rank \(16\) \cite[\S8]{Se}. \(\I_{1,n}\) is the odd unimodular lattice
of signature \((1,n)\) and \(\langle m\rangle\) the rank-one lattice
with generator of square \(m\). A sublattice \(S\subset\Lambda\) is
\emph{primitive} if \(\Lambda/S\) is torsion free; a vector \(k\) is
\emph{characteristic} if \(k\cdot x\equiv x^2\pmod 2\) for all \(x\).
Two facts from \cite{Ni} are used repeatedly: a unimodular sublattice
\(S\subset\Lambda\) is a direct summand, \(\Lambda=S\oplus S^\perp\);
and if \(v\) is a primitive vector of a unimodular lattice with
\(v^2\neq0\), then \(v^\perp\) has cyclic discriminant group of order
\(|v^2|\).

\subsection{The degeneration side}\label{sec:degside}

Let \(\pi\colon\calX\to\Delta\) be as in Theorem~\ref{thm:main}, with
\(C=V_1\cap V_2\) and \(c_i=[C]=-K_{V_i}\in\NS(V_i)\). Put
\begin{equation}\label{eq:limit}
 \xi=(c_1,-c_2)\in\NS(V_1)\oplus\NS(V_2),\qquad
 \Lambda_X=\xi^\perp/\Z\xi=I^\perp/I .
\end{equation}
Here \(I\subset H^2(X_t;\Z)\) is the primitive isotropic plane of the
degeneration (the image of \(\log T\), where \(T\) is the unipotent
monodromy acting on \(H^2(X_t;\Z)\)), and the second equality in \eqref{eq:limit} is an
\emph{integral} isomorphism \cite[Proposition~3.5]{GT}; \(\Lambda_X\)
is even, unimodular and of signature \((1,17)\). The monodromy
invariant lattice \(\calL_t\simeq H\) lies in \(I^\perp\) and meets
\(I\) trivially; we write \(L\subset\Lambda_X\) for its image, and
\begin{equation}\label{eq:SX}
 R_X=L^\perp_{\Lambda_X},\qquad
 S_{X,i}=\im\bigl(c_i^\perp\longrightarrow\Lambda_X\bigr),
 \quad\text{where } c_i^\perp\subset\NS(V_i).
\end{equation}
The maps \(c_i^\perp\to\Lambda_X\), \(D\mapsto[(D,0)]\) or
\([(0,D)]\), are injective, and \(S_{X,1}\perp S_{X,2}\) since
\((D_1,0)\cdot(0,D_2)=0\).

\subsection{The fibration side}\label{sec:fibside}

We paraphrase the framework of \cite[\S\S2,4,5]{GT}, which builds on
the pseudolattice formalism of \cite{Kuz,HT}; the precise definitions
are those of \cite{GT}. Let \(Y\to\PP^1\) be an elliptic K3 surface
with a section and \(\gamma\subset\PP^1\) a loop avoiding the critical
values, with \(\PP^1=\Delta_1\cup_\gamma\Delta_2\) a union of closed
discs, \(Y_i\) the preimage of \(\Delta_i\) and \(F_p\) the fibre over
a base point \(p\in\gamma\). The relative homology \(H_2(Y_i,F_p)\) is
free of rank \(e(Y_i)\), the Euler number of \(Y_i\) (the sum over the
singular fibres in \(\Delta_i\)), generated, when these fibres are nodal,
by Lefschetz thimbles (the discs swept out by the vanishing cycles
along paths from \(p\) to the critical values). It carries the
unimodular \emph{Seifert form} \(\chi_i\) (upper triangular with ones
on the diagonal in a basis of distinguished thimbles, i.e.\ over an
ordered system of disjoint paths, the entries above the diagonal being
intersection numbers of vanishing cycles). The boundary map
\(\phi_i\colon H_2(Y_i,F_p)\to H_1(F_p)\) is the \emph{charge map}, and
its \emph{right adjoint} \(r_i\colon H_1(F_p)\to H_2(Y_i,F_p)\) is
defined by \(\chi_i(v,r_i(x))=\langle\phi_i(v),x\rangle\), where
\(\langle\ ,\ \rangle\) is the intersection form on \(H_1(F_p)\). Fix a
symplectic basis \(a,b\) of \(H_1(F_p)\) with \(\langle a,b\rangle=-1\)
(the sign convention of \cite{GT}). The loop is \emph{allowable}
\cite[Definition~4.11]{GT} if both pieces satisfy
\cite[Assumption~4.9]{GT}: the monodromy around \(\gamma\), traversed
as the anticlockwise boundary of \(\Delta_i\), is
\(\left(\begin{smallmatrix}1&e(Y_i)-12\\0&1\end{smallmatrix}\right)\)
in the basis \((a,b)\), and \(r_i(a)\) is primitive. For such a piece
\(r_i(a)\) plays the role of the point class of a surface-like
pseudolattice \cite[\S\S2,4]{GT}, whose N\'eron--Severi lattice is
\(N_i=r_i(a)^\perp/\Z\,r_i(a)\), where
\(r_i(a)^\perp=\{v:\chi_i(v,r_i(a))=0\}\), with the surface-like form
\(q_i=-\chi_i\) (symmetric there), and whose canonical class is, up to a sign that does
not affect its orthogonal complement, \(k_i=[r_i(b)]\in N_i\). Since
\(a\) is monodromy invariant, transporting it around \(\gamma\) sweeps
out a torus in \(Y\); let \(F\) be the fibre class and
\(\bar\tau\in F^\perp/\Z F\) the class of this torus. The
quasipolarisation \(M\subset\NS(Y)\) is \emph{compatible} with the
fibration and \(\gamma\) \cite[Definition~4.21]{GT} if \(F\in M\) and
\(\bar\tau\) has a lift \(\tau\in F^\perp\subset H^2(Y;\Z)\) with
\(\tau\in M^\perp\). Given such a lift one forms
\begin{equation}\label{eq:LambdaY}
 \Lambda_Y=(\Z F+\Z\tau)^\perp/(\Z F+\Z\tau),\qquad
 \Gamma=F^\perp_M/\Z F\subset\Lambda_Y,\qquad
 S_{Y,i}=\im\bigl(k_i^\perp\longrightarrow\Lambda_Y\bigr),
\end{equation}
where \(k_i^\perp\subset N_i\) and the maps come from the integral
gluing of \cite[Propositions~4.18 and~4.20]{GT}, recalled in
\eqref{eq:glue} below. The four clauses of the \emph{mirror pair} definition
\cite[Definition~5.1]{GT} are verified one by one, each identified by
its number, in \S\ref{sec:final}.

\subsection{Elliptic fibrations}

Singular fibres are named by Kodaira type: \(\I_1\) (nodal, Euler
number \(1\)), \(\II^*\) (Euler number \(10\), components not meeting
the zero section spanning \(E_8\)), \(\I_n^*\) (Euler number \(n+6\),
such components spanning \(D_{n+4}\)). A nodal fibre has monodromy the Picard--Lefschetz
transvection \(T_v\) along its vanishing cycle \(v\in H_1(F_p)\); in
the basis \((a,b)\), with \(\langle a,b\rangle=-1\) and \(v=pa+qb\),
\begin{equation}\label{eq:PL}
 T_{(p,q)}=\left(\begin{smallmatrix}1-pq&p^2\\-q^2&1+pq\end{smallmatrix}\right),
 \qquad\text{so that}\qquad
 T_a=\left(\begin{smallmatrix}1&1\\0&1\end{smallmatrix}\right).
\end{equation}

\section{The restriction period and very generality}\label{sec:period}

The \emph{restriction period} of the central fibre
\cite[Construction~4.3, Proposition~4.4]{AE} is the homomorphism
\(\varphi\colon\Lambda_X\to\Pic^0(C)\) sending the class of
\((D_1,D_2)\in\xi^\perp\) to \(D_1|_C-D_2|_C\). It is well defined:
\(D_1\cdot c_1=D_2\cdot c_2\) on \(\xi^\perp\), so the difference has
degree zero, and \(\varphi(\xi)=c_1|_C+c_2|_C=0\) is exactly
\(d\)-semistability. The inclusion \(L\subset\ker\varphi\) holds
because the classes of \(\calL\) extend to line bundles on all of
\(\calX\), the period arc of \(\calX\) (the image of \(\Delta^*\)
under the marked period map) lying in the \(H\)-locus (the locus of
\(H\)-polarised periods) \cite[\S7E]{AE}, and the
restrictions of a line bundle on \(X_0\) to \(V_1\) and to \(V_2\)
agree on \(C\). Since \(L\simeq H\) is unimodular,
\(\Lambda_X=L\oplus R_X\) with \(R_X\) even, unimodular and negative
definite of rank \(16\). Before using the period domain we lift this
splitting to \(H^2(X_t)\).

\begin{lemma}\label{lem:dec}
There is an orthogonal decomposition
\begin{equation}\label{eq:dec}
 H^2(X_t;\Z)=H_{e_1}\oplus H_{e_2}\oplus L\oplus R_X,\qquad
 H_{e_j}=\Z e_j\oplus\Z e_j^\vee\simeq H,\quad I=\Z e_1\oplus\Z e_2,
\end{equation}
compatible with \eqref{eq:limit}, i.e.\ inducing the identification
\(I^\perp/I=L\oplus R_X\).
\end{lemma}

\begin{proof}
The unimodular \(L\) splits off \(H^2(X_t;\Z)\); the complement
\(L^\perp\) is even unimodular and \(I\) is primitive in it, so
restriction to \(I\) maps \(L^\perp\) onto \(I^\vee\). Choose a basis
\(e_1,e_2\) of \(I\) and \(e_j^\vee\in L^\perp\) with
\(e_i\cdot e_j^\vee=\delta_{ij}\); replace \(e_j^\vee\) by
\(e_j^\vee-\tfrac12(e_j^\vee)^2e_j\) (integral by evenness) and then
\(e_2^\vee\) by \(e_2^\vee-(e_1^\vee\cdot e_2^\vee)e_1\). The new
\(e_j^\vee\) are isotropic and mutually orthogonal, so
\(H_{e_1}\oplus H_{e_2}\) is unimodular, contains \(I\) and splits off;
its complement in \(L^\perp\) maps isomorphically onto
\((I^\perp/I)/L=R_X\).
\end{proof}

\begin{lemma}\label{lem:kernel}
Under the very-generality hypothesis of Theorem~\ref{thm:main},
\(\ker\varphi=L\).
\end{lemma}

\begin{proof}
Because \(\varphi\) vanishes on \(L\), it suffices to show that
\(\varphi(r)\neq0\) for \(0\neq r\in R_X\). Fix \eqref{eq:dec}, in
which \(R_X\) is lifted to a sublattice of \(H^2(X_t;\Z)\) orthogonal
to \(e_1,e_2\).
In a mixed-marked Type~II chart of the period domain near the
boundary point of the degeneration \cite[Theorem~7.1]{AE}, the
coordinates are the modulus \(\tau_C\) of \(C\simeq\C/(\Z+\tau_C\Z)\),
an extension coordinate \(z\in R_X\otimes\C\) and a logarithmic
smoothing coordinate \(w\) (the boundary corresponds to
\(\operatorname{Im}w\to\infty\)). In the decomposition \eqref{eq:dec},
normalised by \((\Omega,e_1)=1\), the period vector is
\(\Omega=e_1^\vee+\tau_Ce_2^\vee+z+we_2-(\tau_Cw+z^2/2)e_1\), the last
coefficient being forced by \(\Omega^2=0\); and the restriction period
\(\varphi(r)\) of \(r\in R_X\) is \((r,z)\bmod(\Z+\tau_C\Z)\), the
extension coordinate recording the period of the central fibre
\cite[Construction~4.3, Proposition~4.13]{AE}.
If this vanishes for some \(r\neq0\), then \((r,z)=m+n\tau_C\) with
\(m,n\in\Z\), and \(r'=r-me_1-ne_2\) is integral, orthogonal to \(L\),
of square \(r'^2=r^2<0\) (as \(I\) is isotropic and orthogonal to
\(r\)), with \((\Omega,r')=(r,z)-m-n\tau_C=0\). So the hyperplane
\(\{(\Omega,r')=0\}\) of periods orthogonal to \(r'\) passes through
the boundary point and meets the interior of the period domain (its
equation does not involve \(w\), so take \(\operatorname{Im}w\)
large): the boundary point lies on the
closure of the Noether--Lefschetz locus of \(r'\), whose image in the
ABE compactification is the closed Noether--Lefschetz locus excluded
by the hypothesis.
\end{proof}

Consequently \(R_X\) is isometric to \(E_8^2\) or to \(D_{16}^+\)
\cite[\S8]{Se}; these are the two Type~II boundary components
(``cusps'') of the \(H\)-polarised period space
\cite[\S7F, Theorem~7.58]{ABE}.

Finally we fix a basis of \(L\). The lattice \(L\simeq H\) contains
exactly two roots \(\pm s\); by
Riemann--Roch one of them is effective on each smooth fibre, and the
effective sign is constant in the smooth family (semicontinuity). Let
\(s\) be the effective root and \(f\) the primitive isotropic class
with \(s\cdot f=1\) lying in the closure of the positive cone of
\(X_t\) (the component containing \(A_t\)); of the two candidates
\(f\) and \(-s-f\) exactly one does. Then
\begin{equation}\label{eq:ab}
 s^2=-2,\qquad f^2=0,\qquad s\cdot f=1,\qquad
 [A_t]=as+bf\quad\text{with } a>0,\ b\geq2a .
\end{equation}
The inequalities hold because \(A_t\) is nef and big:
\(A_t\cdot f=a\geq0\) as both lie in the closed positive cone,
\(A_t\cdot s=b-2a\geq0\) as \(s\) is effective, and
\(A_t^2=2a(b-a)>0\) forces \(a>0\).

\section{Reference models}\label{sec:reference}

The hypotheses give no explicit description of \(X_0\). We therefore
construct, for the given marked period, an explicit Kulikov model
\(\calW\to\Delta\) with the same punctured family; \S\ref{sec:models}
compares \(\calX\) with it. A \emph{surface pair} is a \(d\)-semistable union \(V_1\cup_CV_2\) of
smooth rational surfaces along a common anticanonical elliptic curve,
and \(\Lambda_W\) denotes the lattice \eqref{eq:limit} of such a pair
(it depends only on \(\NS(V_i)\ni c_i\)). An \emph{interior}
\((-2)\)-curve on a component is one disjoint from \(C\).

\begin{proposition}\label{prop:reference}
After shrinking the disc, there is a Kulikov model \(\calW\to\Delta\),
isomorphic over \(\Delta^*\) to the given family together with its
\(H\)-classes and monodromy, without base change, whose central fibre
\(W_1\cup W_2\) is as follows.
If \(R_X=E_8^2\), the \(W_i\) are rational elliptic surfaces
without reducible fibres; the double curve is a fibre \(c_i\) of each,
and there are sections \(s_i\subset W_i\) (\((-1)\)-curves) meeting
\(C\) at the same point, with
\begin{equation}\label{eq:Eref}
 s=(s_1,s_2),\qquad f=(c_1,0)=(0,c_2)\quad\text{in }\Lambda_W .
\end{equation}
If \(R_X=D_{16}^+\), then \(W_1=\F_2\) is the Hirzebruch
surface with negative section \(s_1\), ruling \(g_1\), and
\(c_1=2s_1+4g_1\) (so \(c_1^2=8\) and
\(C\cap s_1=\emptyset\)); \(W_2\) is the blow-up of \(\F_0\) at sixteen
points of a smooth anticanonical curve of bidegree \((2,2)\), lying on
distinct fibres of one ruling \(g_2\); \(W_2\) has no interior
\((-2)\)-curve, and
\begin{equation}\label{eq:Dref}
 s=(s_1,0),\qquad f=(g_1,g_2),\qquad c_2^2=-8 .
\end{equation}
The birational identification \(\calW\dashrightarrow\calX\) is an
isomorphism in codimension one; in particular the divisor \(A\)
itself has a strict transform on \(\calW\), an effective, flat Cartier
divisor.
\end{proposition}

The construction realises the \emph{given} \(\varphi\), not just its
cusp type (compare \cite[Proposition~7.23]{AE}). It uses the
lattices \(\Lambda_W\) of the two pairs,
\(L\oplus E_8^{(1)}\oplus E_8^{(2)}\) (the superscript indicating the
component \(W_i\) carrying the summand) and \(L\oplus D_{16}^+\), with
\(s,f\) as in \eqref{eq:Eref}, \eqref{eq:Dref}, computed in
\S\ref{sec:latticesX} directly from \(\NS(W_i)\) (independently of
Lemma~\ref{lem:models} below). As \(R_X\) is isometric to the rank-\(16\)
summand, we fix an isometry \(\Lambda_W\to\Lambda_X\) sending \(s,f\)
to \(s,f\), and regard \(\varphi\) as a homomorphism on \(\Lambda_W\),
hence on \(E_8^{(i)}\), vanishing on \(L\).

\begin{proof}
\emph{Step 1: a surface pair with period \(\varphi\), case \(E_8^2\).}
Let \(P_i=\Z l_i\oplus\bigoplus_{j=1}^9\Z e_{ij}\) be the Picard
lattice of \(\PP^2\) blown up at nine points, with anticanonical class
\(c_i=3l_i-\sum_je_{ij}\) and \(s_i=e_{i9}\). The plane
\(\langle s_i,c_i\rangle\) is unimodular (Gram matrix
\(\left(\begin{smallmatrix}-1&1\\1&0\end{smallmatrix}\right)\)), so
\(P_i=\langle s_i,c_i\rangle\oplus E_8^{(i)}\). Choose \(P\in\Pic^1(C)\)
and define \(\rho_i\colon P_i\to\Pic(C)\) by \(\rho_i(s_i)=P\),
\(\rho_i(c_i)=0\), \(\rho_1|_{E_8^{(1)}}=\varphi\),
\(\rho_2|_{E_8^{(2)}}=-\varphi\) (the sign comes from the definition
of \(\varphi\)).
Then \(\deg\rho_i(D)=D\cdot c_i\). Embed \(C\) in a plane by the
degree-three bundle \(\rho_i(l_i)\) and blow up the nine points of
\(C\) given by the degree-one bundles \(\rho_i(e_{ij})\). On the
resulting surface \(W_i\) the curve \(C\) is anticanonical with
restriction map \(\rho_i\); as \(\rho_i(c_i)=0\), its normal bundle is
trivial and \(0\to\mathcal O\to\mathcal O(C)\to\mathcal O_C\to0\)
shows that \(|c_i|\) is a basepoint-free pencil with section \(s_i\).
A coincidence among the nine points, or a reducible fibre, would give
a root in \(c_i^\perp\) with zero period (\(e_{ij}-e_{ik}\),
respectively an interior \((-2)\)-curve). Its class in
\(\Lambda_W=L\oplus R_X\) is a root outside \(L\), since the classes
of \(L\) supported on one component are the multiples of the
isotropic \(f\); as \(\varphi\) vanishes on \(L\), its
\(R_X\)-component would be a nonzero element of \(\ker\varphi\),
contradicting Lemma~\ref{lem:kernel}. Both sections meet \(C\) at
\(P\). This constructs \eqref{eq:Eref} with restriction period
\(\varphi\).

\emph{Step 2: case \(D_{16}^+\).}
Choose \(G\in\Pic^2(C)\). The double cover \(C\to\PP^1\) defined by
\(|G|\) embeds \(C\) in \(\F_2\) as an anticanonical curve disjoint
from the negative section, with normal bundle \(4G\). For the other
component consider the abstract lattice
\(P_2=\Z z\oplus\Z g_2\oplus\bigoplus_{j=1}^{16}\Z E_j\) with
\(z^2=g_2^2=0\), \(z\cdot g_2=1\), \(E_j^2=-1\) and
\(c_2=2z+2g_2-\sum_jE_j\),
the Picard lattice of \(\F_0\) blown up at sixteen points, and its
index-two sublattice \(P_2^{\rm ev}=\{D: c_2\cdot D\in2\Z\}\). For
\(D\in P_2^{\rm ev}\) with \(c_2\cdot D=2m\) prescribe
\(\rho_2(D)=mG-\varphi([(mg_1,D)])\), where \((mg_1,D)\in\xi^\perp\)
because \(mg_1\cdot c_1=2m\). This is a homomorphism with
\(\rho_2(g_2)=G\) and \(\rho_2(c_2)=-4G\), since \((g_1,g_2)=f\) and
\((-4g_1,c_2)+\xi=2s\) lie in \(L=\ker\varphi\). Extend \(\rho_2\) to
\(P_2\) by choosing a degree-one square root of \(\rho_2(2E_1)\)
(multiplication by two is surjective on \(\Pic^0(C)\)). Put
\(Z=\rho_2(z)\) and \(p_j=\rho_2(E_j)\in C\). The degree-two bundles
\(Z,G\) are distinct, since otherwise the root \(z-g_2\perp c_2\)
would have zero period; their two maps to \(\PP^1\) embed \(C\) as a
smooth \((2,2)\)-curve in \(\F_0\). Blow up the \(p_j\). They are
distinct and lie on distinct fibres of the ruling \(g_2\), since
otherwise \(E_i-E_j\) or \(g_2-E_i-E_j\) would have zero period; these
negative-square classes lie outside \(L\), whose classes supported on
this component are multiples of \(2s+4f\), of positive square. The
resulting restriction map is \(\rho_2\), and \(\sum_jp_j=2Z+6G\) (the
degree-sixteen part of \(\rho_2(c_2)=-4G\)) says that the normal
bundle of \(C\) in \(W_2\) is \(-4G\): the pair is \(d\)-semistable.
Its restriction period is \(\varphi\) by the defining formula on
\(P_2^{\rm ev}\), and an interior \((-2)\)-curve on \(W_2\) would again
contradict \(\ker\varphi=L\). This is the generic component
\eqref{eq:Dref}, not the special \(\Bl_{16}\F_2\) model of
\cite[\S7F]{ABE}.

\emph{Step 3: a family over the same disc.}
A Tyurin degeneration has monodromy of multiplicity one in the sense
of \cite[Theorem~3.13]{AE}, so its mixed-marked period arc is
transverse to the boundary. Both constructed pairs are
\(d\)-semistable, hence are central fibres of Kulikov models
\cite{Fr83}, and by the Type~II local period theorem
\cite[Theorem~7.1, Corollary~7.24]{AE} the local period map on the
semiuniversal smoothing of such a pair is unramified. Pulling its
local inverse back along the period arc of the \emph{original} family
gives a smooth Kulikov model \(\calW\) over the same disc, with no base
change. The \(H\)-classes extend on the \(H\)-locus \cite[\S7E]{AE},
and the markings can be matched at one smooth point, so the monodromy
generators agree.

\emph{Step 4: identification over \(\Delta^*\).}
For an even integer \(N\gg0\) the class \(A_N=4s+Nf\) has central
extensions \(\bigl(4s_1+\tfrac N2c_1,\ 4s_2+\tfrac N2c_2\bigr)\),
respectively \(\bigl(2s_1+(N-4)g_1,\ Ng_2+c_2\bigr)\), on the two
reference pairs; positivity on every curve follows from the pencils
and the absence of interior roots, so these are ample. They extend on
the \(H\)-locus and remain relatively ample near the reference central
fibre.

On the original smooth fibres \(f\) is nef. Indeed, let
\(d=xs+yf+r\), \(r\perp L\), be an effective root with
\(x=d\cdot f<0\). Nefness of \(A_t\) gives
\(0\leq A_t\cdot d=ay+(b-2a)x\), which with \(x<0\), \(b\geq2a\)
forces \(y\geq0\); then \(-2=d^2=-2x^2+2xy+r^2\) with \(xy\leq0\),
\(r^2\leq0\) forces \(d=-s\), which is not effective. Since \(f\) lies
in the closure of the positive cone and pairs nonnegatively with all
effective roots, it is nef \cite[Ch.~8]{Hu}. For \(N/4>b/a\),
\(A_N=(4/a)A_t+(N-4b/a)f\) is nef on the original fibres, and ample
because an \(A_N\)-orthogonal algebraic root there would be algebraic
on the reference fibre with the same period, where \(A_N\) is ample.
Strong Torelli \cite[Theorem~2.3]{AE} now identifies the two punctured
families with their \(H\)-classes, and the elementary-modification
theorem for Kulikov models \cite[Corollary~3.1]{SB} makes this
identification an isomorphism in codimension one.

Finally, strict transform under a codimension-one isomorphism
preserves horizontal prime divisors, and smoothness of \(\calW\) makes
their sum Cartier and flat. All choices used the marked period of the
given arc, never a coarse boundary point (\S\ref{sec:conventions}).
\end{proof}

\section{Classification of the central fibre}\label{sec:models}

We now determine the central fibre of \(\calX\) itself. The idea is to
transport \(A\) to \(\calW\), make it nef by at most one elementary
modification, and compare the contractions of \(\calW\) and \(\calX\)
defined by the common relative section ring
\(\bigoplus_{m\ge0}\pi_*\mathcal O(mA)\). Componentwise positivity
of \(A\) is what excludes the endpoints of the parameter intervals
below.

\begin{lemma}\label{lem:models}
Up to interchanging the components, the tuple
\((\Lambda_X,L,S_{X,1},S_{X,2})\) of \(\calX\) is that of one of the
following surface pairs:
(\(\mathrm E_0\)) the reference pair \eqref{eq:Eref};
(\(\mathrm E_1\)) the pair obtained from \eqref{eq:Eref} by
contracting \(s_1\) in \(W_1\) and blowing up the point
\(s_1\cap C=s_2\cap C\) on \(W_2\), so that \(c_1^2=1\), \(c_2^2=-1\);
(\(\mathrm D\)) the reference pair \eqref{eq:Dref}, with
\(c_1^2=8\) and \(\NS(W_1)\) even.
\end{lemma}

\begin{proof}
\emph{Step 1: the restrictions of \(A\).}
Transfer the divisor \(A\) to \(\calW\) by
Proposition~\ref{prop:reference}; its restrictions \(D_1,D_2\) to
\(W_1,W_2\) are effective. Two extensions of the line bundle
\(\mathcal O(A)\) from the punctured family differ by a multiple of a
central component, hence \((D_1,D_2)\) is determined by \eqref{eq:ab} up to
adding an \emph{integral} multiple of \(\xi\).

In case \eqref{eq:Eref} this gives, for an integer parameter,
\begin{equation}\label{eq:ED}
 D_1=as_1+xc_1,\qquad D_2=as_2+yc_2,\qquad x+y=b .
\end{equation}
Since \(s_i+c_i\) is nef on \(W_i\) and \((s_i+c_i)\cdot D_i\) equals
\(x\), respectively \(y\), effectivity gives \(x,y\geq0\).

In case \eqref{eq:Dref}, for an integer \(u\),
\begin{equation}\label{eq:DD}
 D_1=(a-2u)s_1+(b-4u)g_1,\qquad D_2=bg_2+uc_2 .
\end{equation}
Intersecting with the nef rulings gives \(D_1\cdot g_1=a-2u\geq0\) and
\(D_2\cdot g_2=2u\geq0\), that is, \(0\leq u\leq a/2\).

\emph{Step 2: a nef model, case \(E_8^2\).}
If \(x,y\geq a\), both divisors in \eqref{eq:ED} are nef, and the
only curves with zero intersection are the sections \(s_1\) (if
\(x=a\)) and \(s_2\) (if \(y=a\)). If \(0\leq x<a\), then \(y>a\), and
we perform the elementary modification along \(s_1\): this
\((-1)\)-curve of \(W_1\) meets \(C\) once, so it is a
\((-1,-1)\)-curve of the threefold; flopping it contracts \(s_1\) in
\(W_1\) and blows up its point of attachment on \(W_2\). Write
\(q\colon W_2'\to W_2\) for this blow-up, \(E\) for its exceptional
curve, \(W_1'\) for the contraction, and primes for strict transforms.
The strict transform of \(A\) restricts to
\begin{equation}\label{eq:flop}
 D_1'=xc_1',\qquad D_2'=q^*(as_2+yc_2)+(x-a)E ,
\end{equation}
the coefficient of \(E\) being forced by
\(D_1'\cdot c_1'=D_2'\cdot c_2'=x\). The pullback of \(c_1'\) to \(W_1\) is \(s_1+c_1\), whose only curve
of zero intersection is \(s_1\); thus \(c_1'\) is ample on \(W_1'\)
and \(D_1'\) is nef, with no curve of zero intersection when \(x>0\).
On \(W_2'\), \(D_2'\cdot E=a-x\), \(D_2'\cdot s_2'=b-2a\) and
\(D_2'\cdot c_2'=x\). For every other irreducible curve \(\Theta\subset W_2\), of
multiplicity \(m\) at the centre, \(c_2\cdot\Theta\geq m\) and
\(s_2\cdot\Theta\geq0\), hence
\(D_2'\cdot\Theta'=a(s_2\cdot\Theta)+y(c_2\cdot\Theta)-(a-x)m\geq
(y-a+x)m\geq0\); if \(c_2\cdot\Theta=0\) then \(\Theta\) is an
irreducible fibre and \(s_2\cdot\Theta=1>0\). So \(D_2'\) is nef, and
for \(0<x<a\) the only possible curve of zero intersection is
\(s_2'\), an interior \((-2)\)-curve, occurring when \(b=2a\). The case
\(y<a\) is symmetric under interchanging the components.

\emph{Step 3: a nef model, case \(D_{16}^+\).}
For \(0\leq u\leq a/2\) both divisors \eqref{eq:DD} are already nef:
on \(\F_2\), a class \(\alpha s_1+\beta g_1\) with \(\alpha\geq0\) is
nef iff \(\beta\geq2\alpha\), and \(b-4u\geq2(a-2u)\) is
\(b\geq2a\); on \(W_2\), every irreducible curve other than \(C\) has
nonnegative intersection with both \(g_2\) and \(c_2\), while
\(D_2\cdot C=2b-8u\geq4a-8u\geq0\). The squares are
\(D_1^2=2(a-2u)(b-a-2u)\) and \(D_2^2=4u(b-2u)\).
For \(0<u<a/2\) the only possible curve of zero intersection is
\(s_1\) (an interior \((-2)\)-curve), when \(b=2a\): indeed
\(D_1\cdot s_1=b-2a\), \(D_1\cdot g_1=a-2u>0\), \(D_2\cdot C>0\), and
each component of a singular fibre of the ruling \(g_2\) has positive
\(c_2\)-degree, since the blown-up points lie on distinct fibres.

\emph{Step 4: the common contraction excludes the endpoints.}
Denote the nef model of Steps 2--3 by \((\calW',D)\); it is isomorphic
to \(\calX\) in codimension one, so by Hartogs extension the relative
section rings of \((\calW',D)\) and \((\calX,A)\) coincide. A
relatively nef extension of a line bundle that is nef and big on the
general fibre is relatively semiample \cite[\S2A]{ABE}, so there are
proper bimeromorphic morphisms
\(\calW'\xrightarrow{\,\mu\,}\calZ\xleftarrow{\,\nu\,}\calX\) to the same
normal relative \(\operatorname{Proj}\). A central component is
contracted to dimension at most one precisely when the restriction of
the nef line bundle to it is not big, and the components of
\(\calW'_0\) and \(X_0\) correspond under the codimension-one
identification. At \(x=0\), \(y=0\), \(u=0\) or \(u=a/2\) one of the
\(D_i\) has square zero, so the corresponding \(V_i\) would be
contracted, contradicting \((A|_{V_i})^2>0\).

\emph{Step 5: resolution counting.}
It remains to see that the remaining nef models determine \(X_0\); we
do this directly, without assuming connectivity of nef models by
divisor-trivial flops. First, \(\nu\) does not contract the double curve
of \(X_0\): the squares of \(C\) on \(V_1,V_2\) sum to zero, so one is
nonnegative, and the Hodge index theorem for the nef and big class
\(A|_{V_i}\) gives \(A|_{V_i}\cdot C>0\) there (a nef and big class
orthogonal to \(C\) would force \(C^2<0\)); the two degrees
\(A|_{V_i}\cdot C\) agree, both being \(\deg A|_C\). Let \(Z_i\) be the normalisations of the
central components of \(\calZ\). By the lists of curves of zero
intersection in Steps 2--3, the minimal resolutions \(U_i\to Z_i\) are
obtained from \(W_i'\) by contracting the \((-1)\)-sections of zero
\(D\)-degree, while interior \((-2)\)-curves of zero degree are
retained. The double curves \(C_i^U\subset U_i\), identified with
\(C\), are smooth with \(K_{U_i}+C_i^U=0\). Each birational morphism
\(V_i\to Z_i\) factors through the minimal resolution \(U_i\), and
\(V_i\to U_i\) is a composition of \(n_i\) point blow-ups. The
identity \(K_{V_i}+C=0\) forces every centre to lie on the successive
strict transform of \(C_i^U\): a blow-up off that curve introduces a
positive discrepancy which later blow-ups cannot cancel. Each such
blow-up lowers \(C^2\) by one, so \(d\)-semistability of \(X_0\) gives
\(n_1+n_2=(C_1^U)^2+(C_2^U)^2\). Moreover, no centre lies over a smooth point of \(\calZ\): a proper
bimeromorphic morphism from a smooth space with trivial canonical
bundle to a normal space is an isomorphism over the smooth locus of
the target. (Over a smooth ball \(U\subset\calZ\), the Jacobian of
\(\nu\) with respect to nowhere vanishing volume forms is a holomorphic
function on \(\nu^{-1}(U)\) vanishing exactly on the exceptional locus
of \(\nu\); it descends through
\(\nu_*\mathcal O_{\calX}=\mathcal O_{\calZ}\) to a holomorphic function
on \(U\), whose zero set is the image of that exceptional locus, of
codimension at least two in \(U\), hence empty.) Note that only proper analytic
contractions and projective surface resolutions have been used;
projectivity of \(\calX\) is not assumed.

Now count. If no \((-1)\)-section was contracted, the right-hand side
of this identity is \((C^2)_{W_1'}+(C^2)_{W_2'}=0\), so \(V_i=U_i=W_i'\): this gives
\(\mathrm E_0\) (for \(x,y>a\)), \(\mathrm E_1\) (for \(0<x<a\) or
\(0<y<a\)) and \(\mathrm D\). If exactly one section was contracted,
say \(s_1\) with \(x=a<y\), the right-hand side is one; away from the
image \(p_0\) of \(s_1\), \(\mu\) is an isomorphism along the double
curve, so \(\calZ\) is smooth there, and the single blow-up is at the
point over \(p_0\), on \(U_1\) (giving back \(W_1\), i.e.\
\(\mathrm E_0\)) or on \(U_2\) (giving \(\mathrm E_1\)). Two sections
are contracted only when \(x=y=a\), so \(b=2a\); they meet \(C\) at
the same point, the count is two, both centres lie over \(p_0\), and a
second blow-up on the same component must be at the unique point of
the strict transform of the double curve over \(p_0\). The
distributions \((1,1),(2,0),(0,2)\) of the two blow-ups give
\(\mathrm E_0\) and the two orders of \(\mathrm E_1\).

\emph{Step 6: the embedded lattices.}
Thus the central pair \((V_1\cup V_2,C)\) is isomorphic to one of the
listed pairs. The construction \eqref{eq:limit}--\eqref{eq:SX} is
intrinsic to the central pair, the isomorphism preserves the
restriction period, and \(L=\ker\varphi\) by Lemma~\ref{lem:kernel};
hence it identifies the \emph{given} embedded \(L\) and both
\(S_{X,i}\) with those of the listed pair, without any choice of a full
cohomology marking.
\end{proof}

\section{The embedded lattices on the degeneration side}\label{sec:latticesX}

We compute the tuple \((\Lambda,L,S_1,S_2)\) of \eqref{eq:limit},
\eqref{eq:SX} for each of the three surface pairs of
Lemma~\ref{lem:models} directly from the pair; by that lemma the
result is the tuple of \(\calX\). In the splitting \(L\oplus R_X\),
with the basis \(s,f\) of \(L\) of \eqref{eq:ab}, one finds the
following; labels \(1,2\) may be interchanged here and below.
\begin{align}
 \mathrm E_0:\quad&
 R_X=E_8^{(1)}\oplus E_8^{(2)},&
 S_{X,i}&=\Z f\oplus E_8^{(i)};                \label{eq:XE0}\\
 \mathrm E_1:\quad&
 R_X=E_8^{(1)}\oplus E_8^{(2)},&
 S_{X,1}&=E_8^{(1)},\qquad
 S_{X,2}=L\oplus E_8^{(2)}.                    \label{eq:XE1}
\end{align}
For the third case write \(h=s+2f\) (so \(h^2=2\)) and
\(\epsilon\colon D_{16}^{+}\to D_{16}^{+}/D_{16}\simeq\Z/2\); then
\begin{equation}\label{eq:XD}
 \mathrm D:\quad R_X=D_{16}^{+},\qquad S_{X,1}=\Z s,\qquad
 S_{X,2}=\{\,mh+r:\ m\in\Z,\ r\in D_{16}^{+},\
                   \epsilon(r)\equiv m\ (\mathrm{mod}\ 2)\,\}.
\end{equation}
\emph{Case \(\mathrm E_0\).} Here
\(\NS(W_i)=\Z s_i\oplus\Z c_i\oplus E_8^{(i)}\) with
\(\left(\begin{smallmatrix}s_i^2&s_ic_i\\s_ic_i&c_i^2\end{smallmatrix}\right)
=\left(\begin{smallmatrix}-1&1\\1&0\end{smallmatrix}\right)\), and
\(c_i^\perp=\Z c_i\oplus E_8^{(i)}\). Passing to \(\xi^\perp/\Z\xi\),
where \((c_1,0)=(0,c_2)=f\) and \((s_1,s_2)=s\), gives \eqref{eq:XE0}.

\emph{Case \(\mathrm E_1\).} Contracting \(s_1\) leaves
\(c_1^\perp=E_8^{(1)}\) unchanged (now \(c_1^2=1\)), so
\(S_{X,1}=E_8^{(1)}\), a unimodular summand of \(\Lambda_X\). On the
other component \(c_2^2=-1\), so \(\Z c_2\) splits off \(\NS(V_2)\) and
\(c_2^\perp\) is unimodular of rank \(10\); its image \(S_{X,2}\) is a
unimodular sublattice of the rank-\(10\) lattice
\(S_{X,1}^\perp=L\oplus E_8^{(2)}\), hence equal to it. This proves
\eqref{eq:XE1}.

\emph{Case \(\mathrm D\).} We first locate the root sublattice
\(D_{16}\subset R_X\). In the basis \(z,g_2,E_j\) of
\(\NS(W_2)=P_2\) from the proof of Proposition~\ref{prop:reference},
the lattice
\begin{equation}\label{eq:D16}
 R_0=c_2^\perp\cap g_2^\perp
 =\Bigl\{-\tfrac12\bigl(\textstyle\sum_jt_j\bigr)g_2+\sum_jt_jE_j:
                     t_j\in\Z,\ \sum_jt_j\equiv0\ (\mathrm{mod}\ 2)\Bigr\}
\end{equation}
has form \(-\sum t_j^2\) on the even-sum vectors \((t_j)\), so it is
\(D_{16}\). An element of \(R_X=L^\perp\) has a representative
\begin{equation}\label{eq:rep}
 \bigl(k(s_1+2g_1),D\bigr),\qquad
 D\cdot g_2=-k,\quad D\cdot c_2=4k ,
\end{equation}
since orthogonality to \(s=(s_1,0)\) forces the first entry to be a
multiple of \(s_1+2g_1\), orthogonality to \(f\) gives
\(D\cdot g_2=-k\), and \((s_1+2g_1)\cdot c_1=4\). Changing the
representative by \(\xi=(2s_1+4g_1,-c_2)\) changes \(k\) by \(2\). The
classes with even \(k\) form \eqref{eq:D16}, and odd \(k\) occurs
(\(k=1\), \(D=-z+3g_2\)), so \(R_X\) is an even unimodular index-two overlattice of \(D_{16}\):
\(R_X=D_{16}^+\), with \(\epsilon\) the parity of \(k\). Next,
\(c_1^\perp=\Z s_1\) in \(\NS(\F_2)\), so \(S_{X,1}=\Z s\). If
\((0,D)\) with \(D\cdot c_2=0\) represents an element of \(S_{X,2}\),
put \(m=D\cdot g_2\); subtracting \(mh=(m(s_1+2g_1),2mg_2)\) leaves the
element \((-m(s_1+2g_1),D-2mg_2)\) of \(R_X\), whose coefficient in
\eqref{eq:rep} is \(k=-m\), so \(\epsilon(r)\equiv m\). Conversely,
given \(mh+r\) with \(\epsilon(r)\equiv m\), choose the representative
of \(r\) with \(k=-m\); then \(mh+r=(0,D+2mg_2)\) lies in \(S_{X,2}\).
This proves \eqref{eq:XD}.

\section{Fibrations and splitting loops on \texorpdfstring{\(Y\)}{Y}}\label{sec:fibrations}

Since \(Y\) is very general, \(\NS(Y)=M\). We use two elliptic
fibrations with section on \(Y\): the \emph{standard} fibration, with
singular fibres \(2\,\II^*+4\,\I_1\), and the \emph{alternate}
fibration, with singular fibres \(\I_{12}^*+6\,\I_1\)
\cite[\S2.2]{DPT}, \cite{CD}.

\paragraph{The standard fibration and the loops \(\gamma_0,\gamma_1\).}
On the base with coordinate \(t\neq0,\infty\), a Weierstrass equation
of the standard fibration can be normalised to the Inose form
\cite{CD}
\begin{equation}\label{eq:Inose}
 z^2=u^3+\alpha u+\beta+t+t^{-1}:
\end{equation}
the \(\II^*\) conditions at \(t=0,\infty\) force the Weierstrass
coefficients to be \(\alpha t^4\) and
\(t^5(\lambda_0+\lambda_1t+\lambda_2t^2)\), and
rescaling \(t,u,z\) gives \eqref{eq:Inose}, with generic parameters
\(\alpha,\beta\). We use two splitting loops, distributing the
singular fibres over the two pieces \(Y_1\mid Y_2\) as
\begin{equation}\label{eq:Eloops}
 \gamma_0\colon\ \II^*+2\,\I_1\ \big|\ \II^*+2\,\I_1,\qquad\qquad
 \gamma_1\colon\ \II^*+\I_1\ \big|\ \II^*+3\,\I_1 .
\end{equation}
Since arbitrary loops with these fibre counts need not have the
monodromy required for allowability, we construct them explicitly.
The standard fibration is the quadratic base change
\(v=\beta+t+t^{-1}\) of the rational elliptic surface
\(z^2=u^3+\alpha u+v\), whose singular fibres are \(\II^*+2\,\I_1\).
Choose the parameters so that some disc containing the two branch
values \(v=\beta\pm2\) contains no critical value of this rational
surface. Over the complement of that disc the cover is trivial, a
disjoint union of two copies of the complement; the boundary of either copy is \(\gamma_0\), with
identity monodromy on each side. Moving one nodal critical value
across this boundary gives \(\gamma_1\): choosing the symplectic basis
so that its vanishing cycle is \(a\), the anticlockwise monodromy on
the side with Euler number \(11\) is \(T_a^{-1}\) and on the other side
it is \(T_a\). These
constructions transport to the given generic parameters: the space of
base changes with distinct branch and critical values is connected,
and the fibrations and loops are topologically locally trivial over it.

\paragraph{The alternate fibration and the loop \(\gamma_D\).}
Taking \(u\) as base coordinate in \eqref{eq:Inose} and \(\eta=tz\)
gives the alternate fibration
\begin{equation}\label{eq:alt}
 \eta^2=t\bigl(t^2+P(u)t+1\bigr),\qquad P(u)=u^3+\alpha u+\beta ,
\end{equation}
with a section, a section of order two (\(t=\eta=0\)), an
\(\I_{12}^*\) fibre at \(u=\infty\) (the Weierstrass invariant \(c_4\)
and the discriminant vanish there to orders \(2\) and \(18\)) and six nodal fibres over
\(P(u)=\pm2\) (the discriminant away from infinity is a nonzero
constant times \(P(u)^2-4\)). We construct a loop \(\gamma_D\) with
the splitting \(4\,\I_1\mid\I_{12}^*+2\,\I_1\). View \eqref{eq:alt} as the pull-back under the cubic \(v=P(u)\) of the
rational elliptic surface \(\eta^2=t(t^2+vt+1)\), with singular fibres
\(\I_4^*+2\,\I_1\). On this rational surface the two nodal vanishing
cycles may be taken to be
\begin{equation}\label{eq:BC}
 B=b,\qquad D=b+2a\qquad(\langle a,b\rangle=-1):
\end{equation}
the product of the two positive nodal Picard--Lefschetz
transformations is the inverse of the \(\I_4^*\) monodromy, of trace
\(-2\); as \(T_BT_D\) has trace \(2-\langle B,D\rangle^2\), the
vanishing cycles have intersection number \(\pm2\), and a choice of
signs and of the symplectic basis gives \eqref{eq:BC}.

Choose the cubic \(P\) so that the inverse image of a disc containing
\(\pm2\) consists of three disjoint discs. A path in the \(v\)-plane
from that disc to a finite branch value of \(P\) lifts to an arc
joining two of the three discs; the two finite branch values of the
generic cubic \(P\) connect in this way two different pairs of discs
(the cover is connected), hence all three. The image of each
connecting arc goes out and back through regular values of the
rational elliptic surface, so parallel transport along it is trivial
on \(H_1(F_p)\). Thus, in a common marking, the vanishing-cycle pairs
in the three discs are \((B,D),(B,D),(B,D)\). A regular
neighbourhood of two of the discs and their connecting arc is a disc;
its boundary is \(\gamma_D\). As before, connectedness of the
complement of the discriminant in the parameter space of cubics
transports this loop to all generic members of \eqref{eq:alt}. In the
convention \eqref{eq:PL},
\begin{equation}\label{eq:monodromy}
 T_BT_D=\left(\begin{smallmatrix}-1&4\\0&-1\end{smallmatrix}\right),
 \qquad
 (T_BT_D)^2=\left(\begin{smallmatrix}1&-8\\0&1\end{smallmatrix}\right)=T_a^{-8},
\end{equation}
which is the monodromy around \(\gamma_D\) on the four-node side, as
required for \(e(Y_1)=4\); the complementary monodromy is its inverse
\(T_a^{8}\), as required for \(e(Y_2)=20\).

\paragraph{Allowability.}
The monodromies computed above are
\(\left(\begin{smallmatrix}1&e(Y_i)-12\\0&1\end{smallmatrix}\right)\)
in all cases: the Euler numbers are \(12,12\) for \(\gamma_0\),
\(11,13\) for \(\gamma_1\) and \(4,20\) for \(\gamma_D\). It remains to
check primitivity of \(r_i(a)\). An integral relative class \(v\) of
charge \(b\), i.e.\ with \(\phi_i(v)=b\), satisfies
\(\chi_i(v,r_i(a))=\langle b,a\rangle=1\), so
\(r_i(a)\) is primitive whenever such a \(v\) exists. For the
splittings \eqref{eq:Eloops} the charge map is surjective already on a
neighbourhood of the \(\II^*\) fibre, since its nodal resolution has
charges generating \(H_1(F_p)\) \cite[Lemma~4.7, Proposition~4.8,
Table~4.1]{GT}, so this applies for every chosen primitive \(a\). For
\(\gamma_D\) each side contains a nodal pair \eqref{eq:BC}, so it
applies directly. Consequently all three loops are allowable in the
sense of \cite[Definition~4.11]{GT}.

\section{Compatibility and the lattices on the fibration side}\label{sec:latticesY}

\paragraph{Compatibility.}\label{sec:compat}
For any of the three loops, the torus class
\(\bar\tau\in F^\perp/\Z F\) of \S\ref{sec:fibside} is primitive and
isotropic by the integral gluing description
\cite[Propositions~4.18 and~4.20]{GT}. Choose a lift \(\tau_0\in F^\perp\), let \(Q\) be the class
of the zero section, and set \(\tau=\tau_0-(\tau_0\cdot Q)F\). The
torus lies over the loop, hence in smooth fibres, so \(\tau_0\) is
orthogonal to \(F\) and to every singular-fibre component, and the
correction makes \(\tau\cdot Q=0\) as well. These classes span
\(\NS(Y)\otimes\Q\) (\(Q\), \(F\) and the fibre components not meeting
\(Q\) span a sublattice of rank \(18=\rk M\) in both fibrations), hence
\(F\in M\) and \(\tau\in M^\perp\). This is compatibility,
including the integral lift required by \cite[Definition~4.21]{GT}.

\paragraph{The frame.}
With \(\Lambda_Y\) and \(\Gamma\) as in \eqref{eq:LambdaY}, since
\(M\) is unimodular and \(Q\cdot F=1\),
\begin{equation}\label{eq:frame}
 \Lambda_Y=H_0\oplus\Gamma,\qquad H_0=\tau^\perp_{M^\perp}/\Z\tau\simeq H,
 \qquad
 \Gamma\simeq E_8^{(1)}\oplus E_8^{(2)}\ \text{(standard)},\quad
 \Gamma\simeq D_{16}^{+}\ \text{(alternate)}.
\end{equation}
Indeed
\(H^2(Y;\Z)=M\oplus M^\perp\) with \(F\in M\), \(\tau\in M^\perp\), so
\(\Lambda_Y=\Gamma\oplus(\tau^\perp_{M^\perp}/\Z\tau)\) orthogonally
(whence \(H_0=\Gamma^\perp\) in \(\Lambda_Y\)), and the second
summand is \(H\) because \(M^\perp\simeq H^{\oplus2}\) and \(\tau\) is
primitive isotropic (primitivity of \(\bar\tau\) implies that of
\(\tau\)). The unimodular plane \(\langle Q,F\rangle\simeq H\) splits
off \(M\), so \(\Gamma\) is even unimodular negative definite of rank
\(16\) containing the fibre-root lattice \(E_8\oplus E_8\), respectively
\(D_{16}\); in the alternate case \(D_{16}\) has index two in
\(\Gamma\), by its determinant \(4\) (or by the two-torsion section).

\paragraph{Pseudolattice facts.}
The gluing results of \cite[Propositions~4.18 and~4.20]{GT}
give the integral identification
\begin{equation}\label{eq:glue}
 \Lambda_Y=(k_1,-k_2)^\perp/\Z(k_1,-k_2)
 \ \subset\ (N_1\oplus N_2)/\Z(k_1,-k_2),
\end{equation}
with \(S_{Y,i}\) the image of \(k_i^\perp\subset N_i\). The two
images are orthogonal, the maps \(k_i^\perp\to\Lambda_Y\) are
injective (no \((0,w)\) with \(w\neq0\) is a multiple of
\((k_1,-k_2)\)), and the classes of fibre components of \(Y_i\) go to
the corresponding fibre roots in \(\Gamma\). For an allowable piece of
Euler number \(n\in\{11,12,13,20\}\) the lattice \(N_i\) is the odd
unimodular lattice \(\I_{1,n-3}\), and \(k_i\) is characteristic,
primitive and of square \(12-n\)
\cite[Theorems~2.18 and~4.10, Lemma~2.20]{GT}; for the pieces with a
\(\II^*\) fibre, primitivity also follows from the surjectivity of
\(\phi_i\) noted in \S\ref{sec:fibrations}. The sign convention for
\(k_i\) does not affect \(k_i^\perp\). As at least one \(k_i\) is
primitive in each case, \((k_1,-k_2)\) spans a saturated subgroup and
\eqref{eq:glue} is the integral construction, not a rational one. The
piece of Euler number \(4\) is computed by hand below.

\paragraph{The splitting \(\gamma_0\).}
Here \(e(Y_i)=12\), so \(k_i\) is primitive, isotropic and
characteristic in \(N_i=\I_{1,9}\). The fibre roots of the \(\II^*\)
fibre on each side give a unimodular \(E_8^{(i)}\subset k_i^\perp\),
which splits off: \(N_i=E_8^{(i)}\oplus U_i\) with \(U_i\ni k_i\)
odd unimodular of signature \((1,1)\). Choose \(\varepsilon_i\in U_i\)
with \(\varepsilon_i\cdot k_i=1\) (unimodularity and primitivity);
since \(k_i\) is characteristic, \(\varepsilon_i^2\) is odd, and adding
a multiple of \(k_i\) we may assume \(\varepsilon_i^2=-1\). Thus
\(N_i=(\Z\varepsilon_i\oplus\Z k_i)\oplus E_8^{(i)}\) and
\(k_i^\perp=\Z k_i\oplus E_8^{(i)}\). In \eqref{eq:glue} the images
of \((k_1,0)\) and \((0,k_2)\) coincide in a primitive isotropic
vector \(\kappa\), and together with the image of
\((\varepsilon_1,\varepsilon_2)\) (square \(-2\), pairing \(1\) with
\(\kappa\)) it spans \(H_0=\Gamma^\perp\). Hence
\(S_{Y,i}=\Z\kappa\oplus E_8^{(i)}\). Mapping \(\kappa\mapsto f\),
\(H_0\to L\) isometrically and the root summands to their namesakes
gives an isometry \(\psi\colon\Lambda_Y\to\Lambda_X\) carrying
\eqref{eq:frame} and the \(S_{Y,i}\) to \eqref{eq:XE0}, with
\(\Gamma\mapsto R_X\).

\paragraph{The splitting \(\gamma_1\).}
Here \(k_1^2=1\) and \(k_2^2=-1\), so \(\Z k_i\) splits off \(N_i\).
The complement \(k_1^\perp\) is unimodular of rank \(8\) and contains
the fibre roots \(E_8^{(1)}\), hence equals them. Writing
\(N_1=\Z k_1\oplus E_8^{(1)}\) and \(N_2=\Z k_2\oplus k_2^\perp\),
the vectors of \((k_1,-k_2)^\perp\) are
\(\alpha(k_1,-k_2)+(e,w)\) with \(e\in E_8^{(1)}\), \(w\in k_2^\perp\),
so \eqref{eq:glue} splits integrally as
\(\Lambda_Y=E_8^{(1)}\oplus k_2^\perp\), with
\(S_{Y,1}=E_8^{(1)}\) and \(S_{Y,2}=k_2^\perp=H_0\oplus E_8^{(2)}\).
The last equality holds because \(k_2^\perp\) is unimodular of
signature \((1,9)\), even (as \(k_2\) is characteristic), and contains
the unimodular fibre roots \(E_8^{(2)}\). This is \eqref{eq:XE1}, and
any isometry \(H_0\to L\) together with the identity on the root
summands gives \(\psi\), again with \(\Gamma\mapsto R_X\).

\paragraph{The splitting \(\gamma_D\).}
Two integral points must be checked: the small side has ``degree
\(8'\)'' rather than \(8\) (its pseudolattice is that of
\(\PP^1\times\PP^1\), with \(N_1\simeq H\) even and
\(k_1^\perp\simeq\langle-2\rangle\), not that of \(\F_1\), with
\(N_1=\I_{1,1}\) and \(k_1^\perp\simeq\langle-8\rangle\)), and the
large side carries the parity condition of \eqref{eq:XD}.

\emph{Small side.} Use distinguished thimbles
\(\theta_1,\dots,\theta_4\) with charges \(B,D,B,D\) as in
\eqref{eq:BC}. In the basis \((a,b)\), with
\(J=\left(\begin{smallmatrix}0&-1\\1&0\end{smallmatrix}\right)\), the
charge matrix, the Seifert matrix and the right adjoint are
\[
 V=\begin{pmatrix}0&2&0&2\\1&1&1&1\end{pmatrix},\qquad
 \chi_1=\begin{pmatrix}1&2&0&2\\0&1&-2&0\\0&0&1&2\\0&0&0&1\end{pmatrix},
 \qquad r_1=\chi_1^{-1}V^{t}J ,
\]
the entries of \(V^tJ\) being \(\langle\phi_1(\theta_i),a\rangle\)
and \(\langle\phi_1(\theta_i),b\rangle\). Direct multiplication gives
\(r_1(a)=(1,-1,-1,1)\) and \(r_1(b)=(-8,6,4,-2)\). Since
\(\chi_1(\theta_i,r_1(a))=1\) for all \(i\), the subgroup
\(r_1(a)^\perp\) consists of the vectors whose coordinates sum to zero. Modulo
\(r_1(a)\) it has basis \(u=\theta_1-\theta_4\),
\(v=\theta_2-\theta_4\), with
\(\bigl(q_1(u,u),q_1(u,v),q_1(v,v)\bigr)=(0,-1,-2)\) for
\(q_1=-\chi_1\), and \(k_1=-4u+2v\). In particular \(N_1\simeq H\) is
even, \(k_1^2=8\), and \(k_1^\perp=\Z v\) with \(v^2=-2\).

Write \(\delta\) for the image of \(v\) in \(\Lambda_Y\). Every
\(D_{16}\) fibre-root class is supported on the other side of the
loop, so is orthogonal to \(\delta\); the saturation of \(D_{16}\) is
\(\Gamma\), so \(\delta\in\Gamma^\perp=H_0\), and it is a primitive
root there. Choose a basis \(\delta,f_0\) of \(H_0\) with \(f_0^2=0\),
\(\delta\cdot f_0=1\), and put \(h_0=\delta+2f_0\). Then
\begin{equation}\label{eq:delta}
 S_{Y,1}=\Z\delta,\qquad \delta^\perp_{\Lambda_Y}=\Z h_0\oplus D_{16}^{+}.
\end{equation}

\emph{Large side.} Here \(N_2=\I_{1,17}\) and \(k_2\) is primitive of
square \(-8\), so \(k_2^\perp\simeq S_{Y,2}\) has determinant \(8\)
and cyclic discriminant group \(\Z/8\). Since
\(S_{Y,2}\perp S_{Y,1}\) and \(\delta^\perp\) has determinant \(2\),
\(S_{Y,2}\) has index two in the right-hand side of \eqref{eq:delta}.
We claim that
\begin{equation}\label{eq:primitiveD}
 S_{Y,2}\cap\Gamma=D_{16}.
\end{equation}
Let \(G_*\) be the relative homology of a neighbourhood of the
\(\I_{12}^*\) fibre and order the two remaining nodal thimbles last,
so \(H_2(Y_2,F_p)=G_*\oplus\Z^2\). The kernel of the charge map on
\(G_*=H_2(U,F_p)\), \(U\) the neighbourhood, is the image of
\(H_2(U)\), that is \(H_2(U)/\Z[F]\), which is spanned by the fibre
components modulo the total fibre: it is exactly the \(D_{16}\)
fibre-root lattice and is primitive in \(G_*\). The two nodal
thimbles have charges \(B,D\) with \(\langle B,D\rangle=2\), as on the
small side, so the last diagonal block of the upper-triangular Seifert
matrix is \(\left(\begin{smallmatrix}1&2\\0&1\end{smallmatrix}\right)\)
and the corresponding right-adjoint equation has right-hand side
\((1,1)^t\); hence the last two coordinates of \(r_2(a)\) are
\((-1,1)\). Hence \(D_{16}\) remains primitive in
\(N_2=r_2(a)^\perp/\Z r_2(a)\): if \(nx=d+m\,r_2(a)\) with
\(d\in D_{16}\), the last two coordinates show \(n\mid m\), and
primitivity in \(G_*\) gives \(x-(m/n)r_2(a)\in D_{16}\). So
\(D_{16}\) is primitive in \(S_{Y,2}\), and its saturation
\(S_{Y,2}\cap\Gamma\) equals \(D_{16}\), proving
\eqref{eq:primitiveD}.

Thus \(S_{Y,2}\) is the kernel of a nonzero character
\(\alpha\colon\Z h_0\oplus D_{16}^{+}\to\Z/2\) whose restriction to
\(D_{16}^+\) is \(\epsilon\), by \eqref{eq:primitiveD}. If
\(\alpha(h_0)=0\), the kernel would be \(\Z h_0\oplus D_{16}\), with
discriminant group \((\Z/2)^3\), not \(\Z/8\). Hence \(\alpha(h_0)=1\)
and
\begin{equation}\label{eq:SY2}
 S_{Y,2}=\{\,mh_0+r:\ \epsilon(r)\equiv m\ (\mathrm{mod}\ 2)\,\}.
\end{equation}
Equations \eqref{eq:frame}, \eqref{eq:delta} and \eqref{eq:SY2} agree
exactly with \eqref{eq:XD}. Map \(\delta,f_0\) to \(s,f\) and identify
the two copies of \(D_{16}^+\) by any isometry (every isometry of
\(D_{16}^+\) preserves its root sublattice \(D_{16}\), hence
\(\epsilon\)). This supplies the integral isometry in the third case.

In all three cases we have constructed an isometry
\(\psi\colon\Lambda_Y\to\Lambda_X\) with
\(\psi(\Gamma)=L^\perp_{\Lambda_X}=R_X\) and \(\psi(S_{Y,i})=S_{X,i}\)
(\(i=1,2\)), for the pairing of loops with central fibres listed in
\S\ref{sec:strategy}.

\section{Lifting and proof of Theorem~\ref{thm:main}}\label{sec:final}

It remains to verify the clauses of \cite[Definition~5.1]{GT}. On the
degeneration side we have \eqref{eq:dec}. On the
fibration side, \(\tau\) is primitive and isotropic in
\(M^\perp\simeq H^{\oplus2}\); in an even unimodular lattice a
primitive isotropic vector has an isotropic dual pairing to \(1\),
spanning with it a hyperbolic plane \(H_\tau\). Likewise \(F\in M\)
with dual \(Q+F\) spans \(H_F\). Hence
\(H^2(Y;\Z)=H_\tau\oplus H_F\oplus H_0\oplus\Gamma\) with
\(M^\perp=H_\tau\oplus H_0\) and \(M=H_F\oplus\Gamma\).
The inclusions \(L\subset\Lambda_X\) and \(\Gamma\subset\Lambda_Y\) are
split, in particular primitive: this is clause (1).
Clause (2) is the componentwise equality \(\psi(S_{Y,i})=S_{X,i}\)
established in \S\ref{sec:latticesY}. For clause
(3), extend \(\psi\) from \(H_0\oplus\Gamma\) by isometries
\(H_\tau\to H_{e_1}\) and \(H_F\to H_{e_2}\) taking \(\tau\) to
\(e_1\) and \(F\) to \(e_2\); this gives
\begin{equation}\label{eq:hatpsi}
 \widehat\psi\colon H^2(Y;\Z)\longrightarrow H^2(X_t;\Z),\qquad
 \widehat\psi(M^\perp)=H_{e_1}\oplus L,
\end{equation}
since \(\psi(H_0)=\psi(\Gamma^\perp)=R_X^\perp=L\). Finally,
\(\tau\) is \(1\)-admissible in the sense of
\cite[Definition~3.30]{GT}: it is primitive and isotropic in
\(M^\perp\), and its isotropic dual \(g\in M^\perp\) satisfies
\((\tau,g)=1\), so \(\tau\) and \(g\) both have divisibility one in
\(M^\perp\). Thus \eqref{eq:hatpsi} is clause (4a) with \(m=1\).

\begin{proof}[Proof of Theorem~\ref{thm:main}]
By Lemma~\ref{lem:models} the degeneration is of type \(\mathrm E_0\),
\(\mathrm E_1\) or \(\mathrm D\), with embedded lattices
\eqref{eq:XE0}, \eqref{eq:XE1} or \eqref{eq:XD}. Choose on the given
very general \(Y\) the standard fibration with \(\gamma_0\), the
standard fibration with \(\gamma_1\), or the alternate fibration with
\(\gamma_D\), respectively. These are allowable
(\S\ref{sec:fibrations}), the \(M\)-quasipolarisation is compatible
(\S\ref{sec:compat}), and \(\psi\), \eqref{eq:hatpsi} verify
\cite[Definition~5.1]{GT}. Hence \(\calX\to\Delta\) and
\((Y\to\PP^1,\gamma)\) form a mirror pair.
\end{proof}

\begin{thebibliography}{ABE22}

\bibitem[ABE22]{ABE}
V.~Alexeev, A.~Brunyate and P.~Engel,
\emph{Compactifications of moduli of elliptic K3 surfaces: stable pair
and toroidal}, Geom.\ Topol.\ \textbf{26} (2022), 3525--3588.

\bibitem[AE23]{AE}
V.~Alexeev and P.~Engel,
\emph{Compact moduli of K3 surfaces},
Ann.\ of Math.\ \textbf{198} (2023), 727--789.

\bibitem[CD07]{CD}
A.~Clingher and C.~F.~Doran,
\emph{Modular invariants for lattice polarized K3 surfaces},
Michigan Math.\ J.\ \textbf{55} (2007), 355--393.

\bibitem[Do96]{Do}
I.~V.~Dolgachev,
\emph{Mirror symmetry for lattice polarized K3 surfaces},
J.\ Math.\ Sci.\ \textbf{81} (1996), 2599--2630.

\bibitem[DHT17]{DHT}
C.~F.~Doran, A.~Harder and A.~Thompson,
\emph{Mirror symmetry, Tyurin degenerations and fibrations on
Calabi--Yau manifolds},
Proc.\ Sympos.\ Pure Math.\ \textbf{96} (2017), 93--131.

\bibitem[DPT25]{DPT}
C.~F.~Doran, J.~Prebble and A.~Thompson,
\emph{Normal forms and Tyurin degenerations of K3 surfaces polarised
by a rank 18 lattice}, Math.\ Nachr.\ \textbf{298} (2025), 806--828.

\bibitem[Fr83]{Fr83}
R.~Friedman,
\emph{Global smoothings of varieties with normal crossings},
Ann.\ of Math.\ \textbf{118} (1983), 75--114.

\bibitem[GT24]{GT}
L.~Giovenzana and A.~Thompson,
\emph{Degenerations and fibrations of K3 surfaces: lattice
polarisations and mirror symmetry}, arXiv:2405.12009v3 (2024).

\bibitem[HT20]{HT}
A.~Harder and A.~Thompson,
\emph{Pseudolattices, del Pezzo surfaces, and Lefschetz fibrations},
Trans.\ Amer.\ Math.\ Soc.\ \textbf{373} (2020), 2071--2104.

\bibitem[Hu16]{Hu}
D.~Huybrechts,
\emph{Lectures on K3 surfaces},
Cambridge Univ.\ Press, 2016.

\bibitem[Ku77]{Ku}
V.~S.~Kulikov,
\emph{Degenerations of K3 surfaces and Enriques surfaces},
Math.\ USSR-Izv.\ \textbf{11} (1977), 957--989.

\bibitem[Kuz17]{Kuz}
A.~Kuznetsov,
\emph{Exceptional collections in surface-like categories},
Sb.\ Math.\ \textbf{208} (2017), 1368--1398.

\bibitem[Ni80]{Ni}
V.~V.~Nikulin,
\emph{Integral symmetric bilinear forms and some of their
applications}, Math.\ USSR-Izv.\ \textbf{14} (1980), 103--167.

\bibitem[Se62]{Se}
J.-P.~Serre,
\emph{Formes bilin\'eaires sym\'etriques enti\`eres \`a discriminant
\(\pm1\)}, S\'em.\ H.~Cartan \textbf{14} (1961--62), exp.\ 14--15.

\bibitem[SB83]{SB}
N.~I.~Shepherd-Barron,
\emph{Extending polarizations on families of K3 surfaces},
in: The Birational Geometry of Degenerations, Progr.\ Math.\
\textbf{29}, Birkh\"auser, 1983, 135--171.

\end{thebibliography}
\end{document}
